\section{Harmonic extensions and Dolbeault compactness}
\label{sec:harmonic-extensions}

The algebraic generators $e_+,e_-$ fixed in the Koszul calculation
of Section~\ref{sec:geometry} determine classes in
$H^1(X_p,L_p^{\pm2})$ over the entire period ball. Write
$\eta_p^\pm$ for their harmonic representatives for $g_p$ and
$h_p^{\pm2}$. Their masses enter the finite bundle energy.
We first identify these coefficients by elliptic problems on
the central geometry, and then prove that no additional
Dolbeault modes concentrate at the nodes.

\begin{proposition}\label{prop:harmonic-masses}
There are closed representatives $\beta_p^\pm$ of $e_\pm$,
supported on a fixed exterior, and sections $f_p^\pm$ such that
\begin{equation}\label{eq:harmonic-section-projection}
 \eta_p^\pm=\beta_p^\pm-\dbar f_p^\pm,
 \qquad \dbar^*\dbar f_p^\pm=\dbar^*\beta_p^\pm.
\end{equation}
The sections have uniform global $L^2$ and $L^\infty$ bounds
and all fixed-exterior derivative bounds. For
$\varepsilon_i\le r<r_0$, where $\varepsilon_i\asymp|\tau_i|^{1/3}$,
\begin{equation}\label{eq:harmonic-cap-mass}
 \int_{\{\varrho_{i,p}<r\}}|\eta_p^\pm|^2\,\dd V_p\le Cr^4.
\end{equation}
The metrics $h_p$, the sections $f_p^\pm$ and the forms
$\eta_p^\pm$ converge in every fixed-exterior derivative
through the full domain $\mathcal B^\circ$ to unique central
limits. In particular
\begin{equation}\label{eq:extension-coefficients}
 c_\pm(p):=\|\eta_p^\pm\|_2^2
        \longrightarrow c_{\pm,0}:=\|\eta_0^\pm\|_2^2>0.
\end{equation}
The full harmonic degree-one endomorphism space is
\begin{equation}\label{eq:harmonic-arrow-form}
 z=(A,B)\longmapsto\eta_p(z)=
   \begin{pmatrix}0&\eta_p^+ A\\ \eta_p^- B&0\end{pmatrix},
\end{equation}
with Hermitian mass
\begin{equation}\label{eq:extension-mass}
 M_p^{\rm ext}(z,z')=c_+(p)\tr(A^*A')+c_-(p)\tr(B^*B').
\end{equation}
\end{proposition}

\begin{proof}
Put $\mathcal F^\pm=\mathcal L^{\pm2}$ on $\cX$.
The cohomology calculation and base change give
$\pi_*\mathcal F^\pm=0$ and $R^1\pi_*\mathcal F^\pm\simeq\OO_{\mathcal B}$.
Leray therefore identifies $H^1(\cX,\mathcal F^\pm)$ with
the sections of the latter line bundle. The fixed frame gives
an actual extension
\begin{equation}\label{eq:total-extension}
 0\longrightarrow\mathcal F^\pm\xrightarrow{j}
 \mathcal U^\pm\xrightarrow{q}\OO_{\cX}\longrightarrow0.
\end{equation}
Its middle term is locally free because the free quotient
locally lifts and splits. The extension restricts to the
chosen class on every fibre, including $X_0$;
see \cite[Tags 01F2 and 0B3A]{Stacks} for these conventions.

Each total node neighbourhood can be chosen Stein, as a
closed hypersurface in a Stein coordinate product. Cartan's
Theorem B \cite[exposé 18, p.~18-7; exposé 19, Section 5]{Cartan}
gives a holomorphic lift $s_i$ of $1$ there.
Choose a global smooth lift $s_{\rm gl}$ on the smooth total
space and disjoint cutoffs $\chi_i$ supported in those
neighbourhoods and equal to one on smaller fixed ones.
Then
\[
 s=s_{\rm gl}+\sum_i\chi_i(s_i-s_{\rm gl}),\qquad
                      \beta=j^{-1}\dbar s
\]
is a smooth splitting and its extension form. The form is
closed and vanishes on the smaller node neighbourhoods.
Its vertical restrictions are the required representatives
$\beta_p$. The fixed exterior is a smooth proper family,
so these representatives have uniform bounds in every
spatial derivative there. This uses the usual smooth
splitting interpretation of extension classes
\cite[Chapter V, Section 14]{Demailly}; no singular
Dolbeault resolution is required.

The geometric comparison used in
\eqref{eq:diameter-ambient} gives radial cutoffs and volume
bounds
\begin{equation}\label{eq:radial-cutoffs}
 |\dd\chi_{r,p}|\le C/r,\qquad
 \Vol\{\varrho_{i,p}<2r\}\le Cr^6,
                         \qquad\varepsilon_i\le r<r_0.
\end{equation}
These follow from bounded-core and conic-end estimates on
the unit Stenzel metric and homothety. The same estimates
hold centrally. Since $u_p$ in
Proposition~\ref{prop:line-he} is uniformly bounded,
subsequences converge on exteriors to bounded solutions of
the central line equation. The normalization has mean zero
in the limit: for a fixed $r$ the exterior integrals
converge, while the omitted integrals are $O(r^6)$.
If two bounded normalized limits differ by $w$, then $w$
is harmonic. A cutoff $\zeta_r$ vanishing at the tips and
equal to one outside radius $2r$ gives
\[
 \int\zeta_r^2|\dd w|^2\le
         4\int|\dd\zeta_r|^2|w|^2\le Cr^4.
\]
Letting $r\downarrow0$ shows that $w$ is constant, and its
mean is zero. Thus the central HE line metric $h_0$ is
unique and $h_p$ converges without subsequences on exteriors.

On either $L_p^{\pm2}$, the operator
$P_p=\dbar^*\dbar$ has a uniform positive section gap by
Proposition~\ref{prop:split-gap}. Solve
$P_pf_p=\dbar^*\beta_p$. The forcing is uniformly bounded
and supported on the fixed exterior. The gap gives a uniform
$L^2$ bound; interior elliptic estimates give all exterior
derivative bounds. The form in
\eqref{eq:harmonic-section-projection} is closed and coclosed,
and hence is the harmonic representative. Orthogonal
projection bounds its norm by $\|\beta_p\|_2$.
Inside each smaller node neighbourhood the forcing is zero,
and the section Bochner identity gives $D_0^*D_0f_p=0$.
Thus
\[
 \Delta|f_p|^2=-2|D_0f_p|^2\le0.
\]
The maximum principle on the smooth cap bounds $|f_p|$
by its uniformly bounded outer boundary values. This
proves its global $L^\infty$ bound. Testing the covariant
harmonic equation with $\chi_{r,p}^2f_p$ gives
\[
 \int\chi_{r,p}^2|D_0f_p|^2
       \le4\int|\dd\chi_{r,p}|^2|f_p|^2\le Cr^4.
\]
Since $\eta_p=-\dbar f_p$ there, this proves
\eqref{eq:harmonic-cap-mass}.

An exterior subsequential limit $f_0$ is bounded and solves
\[
 \dbar_0^*\dbar_0f_0=\dbar_0^*\beta_0
                                      \quad\text{on }X_0^{\rm reg}.
\]
Two such solutions differ by a bounded covariantly harmonic
section. The same global cutoff estimate forces their
difference to be parallel. The exterior class $\pm2D$
used in Proposition~\ref{prop:split-gap} excludes a nonzero
parallel section. This proves uniqueness and full exterior
convergence of $f_p$ and $\eta_p$. Fatou's lemma gives the
central version of \eqref{eq:harmonic-cap-mass}. To compare
norms, split at a fixed radius $r$: exterior integrals
converge and both discarded contributions are $O(r^4)$.
Taking $p\to0$ and then $r\downarrow0$ proves norm convergence.

If $\eta_0$ vanished, $s_0-jf_0$ would be a holomorphic
lift of $1$ on $X_0^{\rm reg}$ in
\eqref{eq:total-extension}. Near a node $s_0$ is holomorphic
and $\beta_0=0$, so $f_0$ is holomorphic there off the node.
Normal Hartogs extension, in a local holomorphic frame,
extends it across the node
\cite[Chapter II, Corollary 7.8]{Demailly}. The lift then
splits the extension on all of $X_0$, contradicting the
nonzero class $e_{\pm,0}$. Thus both limiting masses are
positive. Finally $H^1(\OO)=0$ and
$h^1(L^{\pm2})=1$. The parallel splitting of the endomorphism
bundle therefore gives all harmonic classes in
\eqref{eq:harmonic-arrow-form}. Opposite matrix blocks are
pointwise orthogonal and the trace pairing factors within
each block, proving \eqref{eq:extension-mass}.
\end{proof}

The numbers $c_{\pm,0}$ are specified by the two bounded
central section problems with their algebraically normalized
classes. They are geometric elliptic data, rather than
intersection numbers. Changing a generator rescales its
mass and inversely changes its arrow coordinate. No equality
between the two masses is imposed.

\subsection{A cap estimate for the twisted Dolbeault operator}

For $F_p=L_p^{\pm2}$ or $\OO$, set
\[
 Q_{p,q}(v)=\|\dbar v\|_2^2+\|\dbar^*v\|_2^2,
                      \qquad v\in A^{0,q}(X_p,F_p).
\]
The following estimate controls arbitrary almost-harmonic
forms, beyond the representatives just constructed.

\begin{lemma}\label{lem:dolbeault-cap}
For fixed sufficiently small $s>0$ and all sufficiently
small $p$ depending on $s$, uniformly in $0\le q\le3$,
\begin{equation}\label{eq:dolbeault-cap}
 \|v\|^2_{\{\varrho_i<s\}}
 \le Cs^2Q_{p,q}(v)+C\|v\|^2_{\{s<\varrho_i<2s\}}.
\end{equation}
The constant is independent of $s,p$ and the node-scale
ratios. The same statement holds for split endomorphism
coefficients.
\end{lemma}

\begin{proof}
In the unitary frame associated with a holomorphic line
frame on a total node neighbourhood, the connection form
is $a_p=\tfrac12(\partial-\dbar)\log h_p$ (multiplied by
$\pm2$ for $F_p$). We first claim
\begin{equation}\label{eq:critical-line-smallness}
 \lim_{s\downarrow0}\limsup_{p\to0}
       \sup_{\varrho_i<s}\varrho_i|a_p|=0.
\end{equation}
The reference metric contributes zero, by the ambient lower
bound \eqref{eq:diameter-ambient}. If the claim failed for
the bounded potential $u_p$, there would be points with
$r_j=\varrho_i(x_j)\to0$ and $r_j|\dd u_{p_j}(x_j)|\ge c>0$.
Pass to a subsequence with $\varepsilon_i/r_j\to\theta\ge0$.
The metrics rescaled by $r_j^{-2}$ converge on compact sets
to a complete Stenzel metric if $\theta>0$, and on compact
annuli to the punctured cone if $\theta=0$.
This follows from \eqref{eq:ricci-flat-input}, first for
fixed auxiliary radius and then with that radius decreasing,
and the explicit scaling of the initial cores. The rescaled
bounded scalar forcing tends to zero. Interior $W^{2,t}$
estimates for $t>6$ give a bounded harmonic limit in $C^1$
on compact sets with nonzero derivative at the base point.

A bounded harmonic function on the complete Stenzel metric
is constant: it belongs, by the AC derivative estimates,
to the harmonic space of growth below three in
\cite[Lemma 6.1]{W_Hodge}, whose nonconstant basis elements
have positive growth and independent leading terms.
On the punctured cone, expand in link scalar eigenfunctions.
A coefficient has the form
\[
 b(r)=c_+r^{-2+\sqrt{4+\lambda}}
                         +c_-r^{-2-\sqrt{4+\lambda}}.
\]
Boundedness at zero and infinity leaves only the constant
for $\lambda=0$. Both alternatives contradict the nonzero
derivative. This proves \eqref{eq:critical-line-smallness}.

Use the spin structure defined by the parallel canonical
form on the Calabi--Yau threefold. Its complex spinors are
$\Lambda^{0,*}$ and its Dirac operator is
$\mathcal D=\sqrt2(\dbar+\dbar^*)$.
On the Ricci-flat cone the radial unitary identification gives
\begin{equation}\label{eq:cone-spin-dirac}
 \mathcal D=J\left(\partial_r+\frac{5}{2r}
                                    +\frac{A_\Sigma}{r}\right),
\end{equation}
where $J$ is unitary and $A_\Sigma$ is the self-adjoint link
spin-Dirac operator, with the appropriate chiral sign copies;
see \cite[Section 6, (42)]{JSZ} and
\cite[Sections 2.4 and 3]{LM}.
The five-dimensional Einstein link has scalar curvature
$20$. Friedrich's inequality gives
\[
 |\lambda(A_\Sigma)|\ge
             \sqrt{\frac5{4(5-1)}20}=\frac52;
\]
we use the primary-author formulation in \cite{FK}.
For a compactly supported spinor put $t=\log r$ and
$w=r^3\psi$. Then
\[
 \|\psi\|_2^2=\int\|w\|^2\,\dd t,\qquad
 \|r\mathcal D\psi\|_2^2
 =\int\|(\partial_t+A_\Sigma-\tfrac12)w\|^2\,\dd t.
\]
The cross term integrates to zero. The distance of $1/2$
from the link spectrum is at least two, so
\begin{equation}\label{eq:spin-cone-hardy}
 \|\psi\|_2\le\tfrac12\|r\mathcal D\psi\|_2.
\end{equation}
This is a weighted Dirac estimate, obtained before squaring
the Dolbeault operator.

On the unit Stenzel cap, unit-scale interior estimates on
dyadic annuli give
\[
 \|\rho\nabla\psi\|_2
                  \le C(\|\rho\mathcal D\psi\|_2+\|\psi\|_2).
\]
Its differentiated AC expansion and
\eqref{eq:spin-cone-hardy} absorb the operator difference
on a sufficiently distant end. A cutoff to that end yields
\[
 \|\psi\|_2\le C\|\rho\mathcal D\psi\|_2
                                     +C\|\psi\|_{L^2(K)}
\]
for a fixed compact core $K$. The complete Stenzel Dirac
operator has no $L^2$ kernel: Lichnerowicz's formula
$\mathcal D^2=\nabla^*\nabla$, tested with large-radius
cutoffs, makes an $L^2$ harmonic spinor parallel, and its
infinite-volume constant norm must vanish
\cite[(2.14)]{GL}. If the compact term could not be removed,
unit-norm spinors with $\rho\mathcal D\psi\to0$ would
have a local harmonic $L^2$ limit. It is zero, contradicting
the displayed estimate. We obtain
\begin{equation}\label{eq:stenzel-spin-estimate}
 \|\psi\|_2+\|\rho\nabla\psi\|_2
                       \le C\|\rho\mathcal D\psi\|_2
\end{equation}
for compactly supported spinors on the entire cap.

Choose the auxiliary radius in
\eqref{eq:ricci-flat-input} comparable to, and larger than,
$2s$. On $\varrho<2s$ the initial metric is exactly
Stenzel, and the scaled $C^1$ difference from $g_p$ is at
most $\delta(s)\to0$ for sufficiently small $p$ depending
on $s$. Identify spinors by the positive orthonormal-frame
map. The principal-symbol difference is $O(\delta)$ and
the connection difference is $O(\delta/\varrho)$, so
\[
 \|\varrho(\mathcal D_{g_p}-\mathcal D_{\sigma_p})\psi\|_2
 \le C\delta(s)(\|\varrho\nabla\psi\|_2+\|\psi\|_2).
\]
The line twist adds Clifford multiplication by $a_p$,
whose weighted norm tends to zero by
\eqref{eq:critical-line-smallness}. Homothety of
\eqref{eq:stenzel-spin-estimate} absorbs both errors,
with constants independent of the individual cap scale.
The parallel canonical form supplies the Dolbeault
identification for each metric; it adds no further twist.
Thus $\|\psi\|_2\le C\|\varrho\mathcal D_{F_p}\psi\|_2$
for spinors supported in the actual cap.
Apply this to $\chi_s v$, with $\chi_s=1$ inside radius
$s$ and zero outside $2s$. The commutator is Clifford
multiplication by $\dd\chi_s$, and
$\varrho|\dd\chi_s|\le C$ on its shell.
For a pure $(0,q)$-form,
$\|\mathcal D_{F_p}v\|_2^2=2Q_{p,q}(v)$.
Squaring proves \eqref{eq:dolbeault-cap}. The trivial line
has no connection perturbation; the orthogonal split
endomorphism decomposition proves its assertion.
\end{proof}

\subsection{The central harmonic space and the complementary gap}

The cap estimate has the following compactness consequence.
If $v_j$ is bounded in $L^2$ and $Q_{p_j,q}(v_j)\to0$,
fixed-exterior Dirac estimates give local $H^1$ bounds.
A subsequence converges strongly in local $L^2$ to a smooth
formal harmonic form on $X_0^{\rm reg}$. Fatou's lemma
makes the limit $L^2$. For fixed $s$, pass to the limit
in the shell term of \eqref{eq:dolbeault-cap}; then let
$s\downarrow0$. The shell norm of the central $L^2$ limit
tends to zero, so no mass is lost at any node. Polarization
preserves finite Gram matrices. For eigenforms with
eigenvalue tending to zero, the second-order equation also
gives smooth exterior convergence.

To use this compactness, we identify the central formal
harmonic space. The needed local primitive estimate is
sensitive to the Ricci-flat density. Near a node put
$s=|z|$ for the induced Euclidean radius on $Q(z)=0$ and
$r\asymp s^{2/3}$ for the cone radius. For a line-valued
$(0,1)$-form the bounded line metric and central conical
asymptotics give
\begin{equation}\label{eq:central-densities}
 g_0\asymp s^{-2/3}g_E,
 \qquad \dd V_{g_0}\asymp s^{-2}\dd V_E,
 \qquad |\alpha|_{g_0,h_0}^2\dd V_{g_0}
                \asymp s^{-4/3}|\alpha|_E^2\dd V_E.
\end{equation}
We claim that for fixed nested node neighbourhoods
$U\Subset U'$ and a weakly closed $L^2(g_0,h_0)$ input,
there is a primitive with
\begin{equation}\label{eq:weighted-local-primitive}
 \dbar f=\alpha\quad\text{on }U^{\rm reg},\qquad
 \|r^{-1}f\|_{L^2(g_0,h_0;U)}
                     \le C\|\alpha\|_{L^2(g_0,h_0;U')}.
\end{equation}

Here is the precise weighted application of the local
Koppelman formula. In \cite[Theorem 1.3]{LR} take
$n=3,d=2,\nu=1,p=2,q=1$. The input in
\eqref{eq:central-densities} is Euclidean $L^2$.
At this exponent the punctured maximal domain agrees
with the whole-variety weak domain, as specified in the
source's footnote 1. Its homotopy identity gives
$\alpha=\dbar K\alpha$, with no degree-one projection
term. The coefficient kernels in
\cite[(3.1)--(3.2), Lemma 3.1 and (6.4)--(6.7)]{LR},
with the singular density already included, are bounded by
\[
 C\frac{|z|^\gamma}{|\zeta|^\gamma|\zeta-z|^5},
                              \qquad\gamma\in\{0,1\},
\]
against $\dd V_E$. The desired scalar output density is
$s^{-10/3}|f|^2\dd V_E$. Conjugating input by
$|\zeta|^{-2/3}$ and output by $|z|^{-5/3}$ reduces the
estimate to $L^2(\dd V_E)$ boundedness of
\[
 T_\gamma(z,\zeta)
   =|z|^{\gamma-5/3}|\zeta|^{2/3-\gamma}|\zeta-z|^{-5}.
\]
Use the Schur function $b(x)=|x|^{-2}$.
The integral estimate \cite[Lemma 2.5]{LR} gives
\[
 \int_{U'}|\zeta|^{-a}|\zeta-z|^{-5}\,\dd V_E(\zeta)
                  \le C|z|^{1-a},\qquad1<a<6.
\]
Taking $a=4/3+\gamma$ in one variable and
$a=11/3-\gamma$ in the other proves
\[
 \int T_\gamma(z,\zeta)b(\zeta)\,\dd V_E(\zeta)\le Cb(z),
 \qquad
 \int T_\gamma(z,\zeta)b(z)\,\dd V_E(z)\le Cb(\zeta).
\]
All four exponents $4/3,7/3,11/3,8/3$ lie strictly
between one and six. The Schur test and the homotopy
identity prove \eqref{eq:weighted-local-primitive}.

Let now $\alpha$ be a smooth formal harmonic $L^2$
$(0,1)$-form on $X_0^{\rm reg}$ with values in
$F_0=L_0^2$ or $L_0^{-2}$. Choose a cover with one
small chart for each node and ordinary holomorphic
charts elsewhere. Use \eqref{eq:weighted-local-primitive}
on the node charts and the local Dolbeault lemma on the
others. The differences of the local primitives are
holomorphic on overlaps. They glue the trivial extensions
on these charts to a locally free extension of $\OO_{X_0}$
by $F_0$, hence define a class
\begin{equation}\label{eq:central-cohomology-injection}
 \Xi(\alpha)\in H^1(X_0,F_0).
\end{equation}
Each node lies in just its designated chart, so the
overlap functions are defined away from the nodes and
satisfy the cocycle identity there. A change of local
primitive on a node chart extends holomorphically by
normal Hartogs; changes of primitives or cover therefore
give the same extension class. Coherent GAGA identifies
the analytic cohomology with the algebraically computed
one-dimensional group.

The map \eqref{eq:central-cohomology-injection} is injective.
Indeed, if the extension splits, subtract its local
holomorphic splittings from the primitives. They glue to
$f$ with $\dbar f=\alpha$ globally on the regular locus.
The subtracted holomorphic functions are bounded on smaller
node charts, and $r^{-1}$ is locally square integrable
in conic dimension six. Thus $f/r\in L^2$ near each node
by \eqref{eq:weighted-local-primitive}. For the global
cutoff $\zeta_r$ used above,
\[
 \dbar(\zeta_r f)=\zeta_r\alpha+(\dbar\zeta_r)f
                           \longrightarrow\alpha\quad\text{in }L^2.
\]
The second term tends to zero by absolute continuity of
the integrals of $|f/r|^2$ on shrinking shells.
Formal coclosedness, tested against the compactly supported
$\zeta_r f$, now gives $\|\alpha\|_2^2=0$.
By Proposition~\ref{prop:harmonic-masses}, the central
space contains the nonzero form $\eta_0^\pm$.
It consequently has dimension exactly one. This argument
uses formal harmonic forms on the regular locus and does
not presume an identification of central Laplacian domains.

\begin{proposition}\label{prop:full-dolbeault-gap}
The split endomorphism Dolbeault Laplacian has a uniform
positive lower bound on the orthogonal complement of its
harmonic kernel, in every degree $0\le q\le3$. Its complex
harmonic dimensions in those degrees are $8,8,8,8$.
\end{proposition}

\begin{proof}
If positive degree-one eigenvalues for $L_p^{\pm2}$ tended
to zero, normalized eigenforms orthogonal to
$\eta_p^\pm$ would, together with
$\eta_p^\pm/\|\eta_p^\pm\|_2$, converge to two orthonormal
central formal harmonic forms. The preceding compactness
preserves both norms and their pairing. This contradicts
the one-dimensional central space. The degree-zero gap is
Proposition~\ref{prop:split-gap}.

On each compact smooth fibre the Hodge--Serre operator is
a conjugate-linear isometry from $F$-valued $(p,q)$-forms
to $F^*$-valued $(3-p,3-q)$-forms and satisfies
\[
 \Delta''_{F^*}(\#v)=\#\Delta''_Fv;
\]
see \cite[Chapter VI, Section 3.1 and proof of Theorem 7.3]{Demailly}.
For $p=0$ the parallel canonical form identifies its target
with $(0,3-q)$-forms valued in $F^*$. The dual metric of
$h_p^2$ is $h_p^{-2}$. Thus the degree-one gap for the
dual line gives degree two, and degree zero gives degree
three, including multiplicities. The harmonic dimensions
for each nontrivial line are $0,1,1,0$.

For the trivial line, the scalar and degree-one gaps
follow from the Hodge input and $\Delta_d=2\Delta''$;
$H^1(X_p;\CC)=0$. Hodge--Serre duality gives the remaining
degrees. The harmonic dimensions are $1,0,0,1$.
Finally
\[
 \End E^{\rm sp}_p
 =(\OO\otimes\End\CC^2)^{\oplus2}
 \oplus(L_p^2\otimes\End\CC^2)
 \oplus(L_p^{-2}\otimes\End\CC^2)
\]
is an orthogonal parallel decomposition. Summing the
dimensions gives $8,8,8,8$ and taking the minimum of the
component gap constants proves the result.
\end{proof}
